\documentclass[11pt]{article}
\usepackage[utf8]{inputenc}
\usepackage[T1]{fontenc}
\usepackage{geometry}
\geometry{a4paper, margin=1in}
\usepackage{charter} % Professional serif font
\usepackage{xcolor}
\usepackage{titlesec}
\usepackage{enumitem}
\usepackage{hyperref}

\definecolor{teal}{HTML}{0D9488}
\definecolor{slate}{HTML}{334155}

\titleformat{\section}{\Large\bfseries\color{slate}\sffamily}{\thesection}{1em}{}[\titlerule]
\titleformat{\subsection}{\large\bfseries\color{teal}\sffamily}{\thesubsection}{1em}{}

\hypersetup{
    colorlinks=true,
    linkcolor=teal,
    urlcolor=teal,
}

\begin{document}

\begin{center}
    {\huge \bfseries \color{slate} Behavioral Digital Twin (BDT) Framework} \\
    \vspace{0.2cm}
    {\large \color{teal} Architectural Blueprint \& Strategic Vision} \\
    \vspace{0.4cm}
    {\small \color{gray} Author: Yichao Jin, Ph.D. \textbullet \ Today: \today}
\end{center}

\vspace{0.8cm}

\section*{Core Definition}
The \textbf{Behavioral Digital Twin (BDT) Framework} is an integrated computational behavioral science system that merges Discrete Choice Experiments (DCE), structural modeling, and generative simulation. It transforms static human behavioral observations into high-fidelity, interactive \textbf{Digital Agents}, enabling \textbf{\textit{in-silico}} testing of policy interventions, product designs, or clinical protocols.

\section*{The Three-Tier Architecture}

\subsection*{1. Neural-Behavioral Ingestion Layer}
\begin{itemize}[leftmargin=*]
    \item \textbf{Function:} Aggregates multi-modal data including experimental results, digital footprints, and psychometric scales.
    \item \textbf{Detail:} Uses non-linear dimensionality reduction and feature alignment to capture latent behavioral parameters (e.g., risk preference, time discount rates, cognitive load), constructing a precise \textbf{Behavioral DNA} for each agent.
\end{itemize}

\subsection*{2. Probabilistic Decision Computing Engine}
\begin{itemize}[leftmargin=*]
    \item \textbf{Function:} A theory-agnostic core that can embed structural econometric models (RUM/RMT) or deep-learning decision networks.
    \item \textbf{Detail:} Leverages Bayesian inference or Inverse Reinforcement Learning (IRL) to reconstruct the "decision logic" within the digital twin. It enables the simulation of dynamic responses to \textbf{counterfactual} scenarios rather than simple outcome prediction.
\end{itemize}

\subsection*{3. Synthetic Simulation \& Intervention Layer}
\begin{itemize}[leftmargin=*]
    \item \textbf{Function:} Deploys thousands of BDT agents into controlled virtual societies or market environments for stress testing.
    \item \textbf{Detail:} Generates massive \textbf{Synthetic Micro-data} through Monte Carlo simulations. This allows researchers to assess marginal effects and identify "vulnerable" or "resistant" sub-populations before any real-world resource allocation.
\end{itemize}

\section*{Key Value Propositions}

\begin{itemize}[leftmargin=*]
    \item \textbf{From Description to Synthesis:} Shifts focus from "what happened" to "what would happen if X changed."
    \item \textbf{Low-Cost "Pre-Flight" Testing:} Increases the success rate of real-world RCTs by over 40\% through pre-screening interventions in a digital sandbox.
    \item \textbf{Interdisciplinary Compatibility:}
    \begin{itemize}
        \item \textbf{Public Health:} Predicting vaccine hesitancy diffusion.
        \item \textbf{Neuroeconomics:} Modeling how cognitive decay impacts long-term financial stability.
        \item \textbf{Policy Science:} Evaluating the equity impacts of subsidy redistribution.
    \end{itemize}
\end{itemize}

\vspace{1cm}
\begin{center}
    \small \color{gray} \textit{For inquiries or collaboration, please visit \href{https://yichao2022.github.io}{yichao2022.github.io}}
\end{center}

\end{document}
